\section{Finite shape bounds}\label{sec:convex}

This section proves that the second difference of $P_N$ changes sign at
most four times. At every even crossing with $N\ge8$ this gives the unique
interior minimum and the two intervals of modes, with no asymptotic input.

For the fair-start second differences, $0<\varepsilon<\frac12$ and $c=1-2\varepsilon$, so
that $r_t(b)=\varepsilon+cb/t$. For $t\ge1$ and every integer $a$ we write
$P_t(a)=\Prob(A_t=a)$, which is $0$ unless $0\le a\le t$, and
\[
 D_t(a)=P_t(a+1)-2P_t(a)+P_t(a-1).
\]
We also put $s_t=D_t(1)=P_t(0)-2P_t(1)+P_t(2)$ for $t\ge2$.

\begin{lemma}[recurrence for the second difference]\label{lem:D-rec}
For every $t\ge1$ and every integer $a$,
\begin{align}
 P_{t+1}(a)&=\bigl(1-r_t(a)\bigr)P_t(a)+r_t(a-1)\,P_t(a-1),\label{eq:P-rec}\\
 D_{t+1}(a)&=\bigl(1-r_t(a+1)\bigr)D_t(a)+r_t(a-2)\,D_t(a-1).\label{eq:D-rec}
\end{align}
Also $P_t(a)=P_t(t-a)$ and $D_t(a)=D_t(t-a)$. The same two recurrences hold
with $f_t(m)$ in place of $P_t(a)$.
\end{lemma}

\begin{proof}
Given $A_t=a$, the copied step is $+1$ with probability $a/t$, so step $t+1$
is $+1$ with probability $(1-\varepsilon)\frac at+\varepsilon(1-\frac
at)=r_t(a)$. This gives \eqref{eq:P-rec} for $0\le a\le t+1$, where at $a=0$ and at
$a=t+1$ the formula has one extra term, namely $r_t(-1)P_t(-1)$ or
$(1-r_t(t+1))P_t(t+1)$, and this term is $0$ because $P_t(-1)=P_t(t+1)=0$. For
other $a$ every term vanishes. The
symmetry of $P_t$ is Lemma~\ref{lem:erw-reduction}, and it gives the symmetry
of $D_t$. The chain $M$ of Lemma~\ref{lem:erw-reduction} has the same
transition probabilities, so \eqref{eq:P-rec} holds for $f_t$ as well.

For \eqref{eq:D-rec} put $\varsigma=c/t$, so that $r_t(b+1)=r_t(b)+\varsigma$.
We insert \eqref{eq:P-rec} at $a+1$, $a$ and $a-1$ into the definition of
$D_{t+1}(a)$ and collect the coefficients of $P_t(a+1)$, $P_t(a)$, $P_t(a-1)$
and $P_t(a-2)$. They are
\[
 1-r_t(a+1),\qquad 3r_t(a)-2,\qquad 1-3r_t(a-1),\qquad r_t(a-2).
\]
The right side of \eqref{eq:D-rec} has the coefficients $1-r_t(a+1)$,
$-2+2r_t(a+1)+r_t(a-2)$, $1-r_t(a+1)-2r_t(a-2)$ and $r_t(a-2)$. The second
ones agree because $2r_t(a+1)+r_t(a-2)=3r_t(a)$, and the third ones agree
because $r_t(a+1)+2r_t(a-2)=3r_t(a-1)$.
\end{proof}

Trinas and Valle~\cite[(2.2)]{tv2026} use \eqref{eq:P-rec}, the forward
equation of Lemma~\ref{lem:erw-reduction}, for a superdiffusive local
limit theorem. Their higher-difference recurrence~\cite[(4.10)]{tv2026}
supports a certified fixed-memory study of log-concavity. We use
\eqref{eq:D-rec} for signs: its coefficients are positive for
$2\le a\le t-1$, so sign changes enter only through the two edges.
The next lemma gives the required variation diminishing property.

For a finite list $z$ of real numbers let $\operatorname{sp}(z)$ be its sign
pattern. We delete the zeros of $z$, replace each entry by its sign, and merge
equal neighbors. So $\operatorname{sp}(z)$ is a list of alternating signs. For
two such lists we write $p\preceq q$ if $p$ is a subsequence of $q$. We put
$\Pi_5=(+,-,+,-,+)$ and $\Pi_3=(+,-,+)$.

\begin{lemma}[one step]\label{lem:interleave}
Let $\zeta=(\zeta_1,\dots,\zeta_n)$ be real, and let
$\eta_i=A_i\zeta_i+B_i\zeta_{i+1}$ for $1\le i\le n-1$ with $A_i\ge0$ and
$B_i\ge0$. Let $w=(\zeta_1,\eta_1,\zeta_2,\eta_2,\dots,\eta_{n-1},\zeta_n)$.
Then $\operatorname{sp}(w)=\operatorname{sp}(\zeta)$ and
$\operatorname{sp}(\eta)\preceq\operatorname{sp}(\zeta)$.
\end{lemma}

We prove this in Appendix~\ref{app:signs-step}.

\begin{lemma}[the sequence next to the end]\label{lem:level-one}
For every $t_1\ge2$ we have $\operatorname{sp}(s_2,s_3,\dots,s_{t_1})\preceq\Pi_3$.
Also, for every $0\le\varepsilon\le1$, the numbers $s_3$ and $s_4$ do not both vanish.
\end{lemma}

We prove this in Appendix~\ref{app:signs-level}.

\begin{theorem}[at most five blocks]\label{thm:five-blocks}
Let $N\ge3$ and $0<\varepsilon<\frac12$. Then
\[
 \operatorname{sp}\bigl(D_N(1),D_N(2),\dots,D_N(N-1)\bigr)\preceq(+,-,+,-,+).
\]
So the second difference of $P_N$ changes sign at most four times on
$\{1,\dots,N-1\}$. The same holds for the second difference of $f_N$, and for
$f_N$ it holds for $N=2$ as well.
\end{theorem}

We prove this in Appendix~\ref{app:signs-five}.

A negative-to-positive change at $a_0<a_1$ and its reflection give
$-,+,+,-$ at $a_0<a_1<N-a_1<N-a_0$. A negative second difference
between them would give the fifth alternating sign forbidden by
Theorem~\ref{thm:five-blocks}. For a zero, trace \eqref{eq:D-rec}
backwards: positivity of its coefficients forces two zero parents,
unless the parents have opposite signs, again giving five alternating
signs. Reaching $t=5$ forces $s_3=s_4=0$, contrary to
Lemma~\ref{lem:level-one}. This proves a whole convex stretch.

\begin{theorem}[no holes]\label{thm:no-holes}
Let $N\ge5$ and $0<\varepsilon<\frac12$. Let $1\le a_0<a_1\le N/2$ be integers
with $D_N(a_0)<0<D_N(a_1)$. Then $D_N(a)>0$ for every $a_1\le a\le N-a_1$.
\end{theorem}

We prove this in Appendix~\ref{app:signs-holes}.

\begin{lemma}[the center from no holes]\label{lem:center-no-holes}
Let $N\ge6$ be even and $0<\varepsilon<\frac12$. If
$P_N(1)>\max\{E_N,C_N\}$, then
$P_N(a)>C_N$ for $1\le a\le N-1$, $a\ne h$, and $D_N(h)>0$.
\end{lemma}

\begin{proof}
Suppose that $P_N(a)\le C_N$ for some $1\le a<h$, and let $k$ be the first
minimizer of $P_N$ on $\{1,\dots,h\}$. Then $2\le k<h$ and
$P_N(k-1)>P_N(k)\le P_N(k+1)$, so $D_N(k)>0$.
The first maximizer $a_0$ on $\{0,\dots,k\}$ has value at least $P_N(1)$,
which exceeds both endpoint values of this interval. Thus $1\le a_0<k$
and $D_N(a_0)<0$. Theorem~\ref{thm:no-holes} gives $D_N(a)>0$ for
$k\le a\le N-k$.

Put $\delta(a)=P_N(a+1)-P_N(a)$. These first differences are strictly
increasing on $\{k-1,\dots,N-k\}$, and symmetry gives
$\delta(h)=-\delta(h-1)$. Hence $\delta(h-1)<0$, and every difference
from $k$ to $h-1$ is negative. This makes $P_N(k)>P_N(h)$, a contradiction.
Reflection proves the other half, and
$D_N(h)=2(P_N(h-1)-P_N(h))>0$.
\end{proof}

\begin{corollary}[the shape at every even crossing]\label{cor:unique-min}
For every even $N\ge8$, at the crossing,
\[
 P_N(a)>C_N=E_N\quad(1\le a\le N-1,\ a\ne h),\qquad D_N(h)>0.
\]
There are integers $1\le l\le r<h$ such that $P_N$ strictly increases
from $0$ to $l$, is constant from $l$ to $r$, and strictly decreases
from $r$ to $h$. Its two intervals of global maximizers are
$\{l,\dots,r\}$ and $\{N-r,\dots,N-l\}$.
\end{corollary}

\begin{proof}
At the crossing $P_N(1)>E_N=C_N$ by Proposition~\ref{prop:dip-hyp}, so
Lemma~\ref{lem:center-no-holes} gives the minimum and central curvature.
For the humps, let $p=\operatorname{sp}(D_N(1),\dots,D_N(h))$.
It ends in $+$, and symmetry makes the full pattern $p$ followed by its
reverse, with the central positive blocks merged. If $p$ has $j$ blocks,
the full pattern has $2j-1\le5$ by Theorem~\ref{thm:five-blocks}.
Thus $p$ is a subsequence of $(+,-,+)$.

The first differences $\delta(a)=P_N(a+1)-P_N(a)$ on $0\le a<h$
therefore have at most three successive monotonic stretches, namely
nondecreasing, nonincreasing, nondecreasing. Their first value is positive,
and their last is $-D_N(h)/2<0$. The first stretch stays positive and the
last stays negative. The middle stretch passes through zero at most once,
possibly with consecutive zeros. This gives the claimed rise, plateau
and fall. Symmetry gives the other hump.
\end{proof}

\begin{corollary}[a window from the maximum]\label{cor:height-window}
Let $N\ge8$ be even and $x=x_N$. Let $[l,r]$ be the left
interval of maximizers in Corollary~\ref{cor:unique-min}, and put
\[
 M=P_N(r),\qquad B=r-1+M^{-1},\qquad
 b=\min\{h,\lceil B\rceil-1\}.
\]
Then $D_N(a)>0$ for every integer $b\le a\le N-b$.
This window is nonempty.
\end{corollary}

\begin{proof}
Let $q$ be the first integer after $l$ with $D_N(q)>0$.
Corollary~\ref{cor:unique-min} gives $D_N(l)<0<D_N(h)$,
so $q\le h$. The constant plateau gives $q>r$.
By its definition $P_N$ is concave on $[r,q]$. Hence
\[
 \frac{M(q-r+1)}2
 \le\sum_{a=r}^{q-1}P_N(a)
 \le\frac{1-P_N(h)}2<\frac12,
\]
where the first inequality sums the chord from $P_N(r)=M$
to the nonnegative value $P_N(q)$. Thus $q<B$, and consequently
$q\le b\le h$. Theorem~\ref{thm:no-holes}, applied at $l<q$,
gives positivity on $[q,N-q]$, which contains $[b,N-b]$.
\end{proof}

Both corollaries use the finite endpoint-dip certificates of
Proposition~\ref{prop:dip-hyp}, but no mode-location or convex-window
estimate. The height bound supplies a window at every even length,
possibly just the center. Heuristically, Theorem~\ref{thm:landau}(b) and
Corollary~\ref{cor:mode-location} give
$M=P_N(r)=\frac{f_L(\lambda_{\mathrm L})}{2x}(1+o(1))$, hence, with $m^*=r$,
$B=m^*+\bigl(2/f_L(\lambda_{\mathrm L})+o(1)\bigr)x$. Since
$f_L(\lambda_{\mathrm L})\approx0.1807$, this is about $m^*+11x$, far below
$h$ for large $N$; at $N=4000$, computed $B-m^*$ is about $9.9x$.
The asymptotic expression is not an explicit finite-length estimate.
The height-based window itself needs no length threshold.
